\begin{tikzpicture}
\begin{semilogxaxis}[
    xlabel={Frecuencia $\omega$ (rad/s)},
    ylabel={Fase (grados)},
    grid=both,
    grid style={line width=.1pt, draw=gray!30},
    major grid style={line width=.2pt, draw=gray!60},
    xmin=1e3, xmax=1e8, % Amplié un poco el rango para ver el final en -180°
    ymin=-200, ymax=20,
    ytick={0, -45, -90, -135, -180}, % Marcas típicas de fase
    width=12cm,
    height=8cm,
    legend style={
    nodes={scale=0.7, transform shape} % Escala todo al 80%
}
]

% Trazado asintótico exacto de fase
% (1e3, 0) -> Inicio plano
% (3691.95, 0) -> 0.1*wp1: Empieza a caer el primer polo
% (49621.45, -50.78) -> 0.1*wp2: Empieza a caer el segundo polo (se suma a la caída del primero)
% (369195.2, -129.22) -> 10*wp1: El primer polo se clava en -90° (solo sigue cayendo el segundo)
% (4962144.8, -180) -> 10*wp2: Ambos polos terminan de caer.
% (1e8, -180) -> Final plano

\addplot[thick, blue] coordinates {
    (1e3, 0)
    (3691.95, 0)
    (49621.45, -50.78)
    (369195.2, -129.22)
    (4962144.8, -180)
    (1e8, -180)
};
\addlegendentry{Teórico (Asintótico)}

% Marcas para indicar dónde empieza el efecto de cada polo
\addplot[mark=|, mark size=3pt, red, forget plot] coordinates {(3691.95, 0)};
\node[anchor=south west, font=\small, text=black] at (axis cs:3691.95, 0) {$0.1\omega_{p1}$};

\addplot[mark=|, mark size=3pt, red, forget plot] coordinates {(49621.45, -50.78)};
\node[anchor=south west, font=\small, text=black] at (axis cs:49621.45, -50.78) {$0.1\omega_{p2}$};

\addplot[thick, dashed, orange] table [
    skip first n=1,
    x expr=\thisrowno{0} * 2 * pi, % Convertimos Hz a rad/s
    y index=2 % Tomamos los grados directamente
] {sources/bode.txt};
\addlegendentry{Simulación}

\end{semilogxaxis}
\end{tikzpicture}
